\chapter{Problem Statement}
\label{chap:problem-statement}

% ============================================================
\section{Motivation}
\label{sec:motivation}
% ============================================================

The clinical deployment of deep learning systems for medical image analysis is
fundamentally constrained by the fragmentation of patient data across healthcare
institutions.
Regulatory frameworks governing the processing of health records --- most notably
the General Data Protection Regulation in the European Union
\cite{gdpr2016regulation} and the Health Insurance Portability and Accountability
Act in the United States \cite{hipaa1996health} --- impose strict conditions on any
cross-institutional transfer of protected health information, conditions that a
conventional centralised training pipeline cannot satisfy without substantial legal
and administrative overhead.
As a consequence, the large, diverse, and representative datasets that deep
learning models require for reliable generalisation are, in practice, inaccessible
through centralisation: each institution's imaging archive constitutes an isolated
silo that can neither be pooled with nor directly compared against those of
neighbouring sites.

Breast cancer molecular subtype classification from dynamic contrast-enhanced MRI
exemplifies this tension.
Accurately distinguishing Luminal~A, Luminal~B, HER2-enriched, and Triple-Negative
subtypes from volumetric DCE-MRI data is clinically consequential --- subtype
assignment directly determines therapeutic strategy --- yet the inter-institutional
variability in scanner hardware, acquisition protocol, patient demographics, and
subtype prevalence means that models trained at a single site are unlikely to
generalise to other clinical environments \cite{ghosh2020clustered}.
Federated learning, which trains a shared model by exchanging only parameter
updates rather than raw patient data, offers a principled mechanism for exploiting
multi-institutional data under these constraints.
To the best of the authors' knowledge, no prior published work has applied federated
learning to molecular subtype classification on the publicly available breast
DCE-MRI cohort \cite{alves2023breastmri} with a variance-reduction aggregation
strategy such as SCAFFOLD, leaving a significant gap in the empirical literature
that this thesis addresses.

% ============================================================
\section{Research Questions}
\label{sec:research-questions}
% ============================================================

This thesis is organised around three research questions that progress from
centralised baseline characterisation to federated optimisation:

\begin{enumerate}

  \item[\textbf{RQ1.}] What classification performance ceiling does a centralised
  deep learning system achieve on breast DCE-MRI molecular subtype prediction
  under severe class imbalance, and how does a clinically motivated binary
  reformulation address the data scarcity constraint imposed by rare subtypes?

  \item[\textbf{RQ2.}] How much classification performance is lost when training
  is distributed across three simulated hospital clients with non-IID data
  partitioning under the FedAvg aggregation strategy, and what is the
  quantitative privacy-performance trade-off?

  \item[\textbf{RQ3.}] Can SCAFFOLD's control-variate correction mechanism reduce
  the convergence gap caused by statistical heterogeneity in a three-hospital
  federated breast MRI setting, and to what extent does convergence-aware
  aggregation recover performance lost to client drift?

\end{enumerate}

% ============================================================
\section{Formal Hypotheses}
\label{sec:hypotheses}
% ============================================================

The three research questions yield the following testable hypotheses:

\begin{enumerate}

  \item[\textbf{H1.}] A ConvNeXt-Nano backbone with Gated Attention MIL, trained
  with LDAM loss and ASAM optimisation, achieves higher macro F1 on the four-class
  centralised task than a standard cross-entropy baseline, because the combination
  of margin-based loss and sharpness-aware minimisation better handles the
  10.4:1 class imbalance present in the dataset.

  \item[\textbf{H2.}] FedAvg aggregation across three non-IID hospital clients
  incurs a performance penalty relative to the centralised baseline, because local
  gradient updates on heterogeneous data distributions cause client drift that is
  not corrected by simple weighted averaging.

  \item[\textbf{H3.}] SCAFFOLD's control-variate mechanism reduces the client
  drift identified in~H2 and narrows the centralised-federated performance gap,
  because it explicitly estimates and corrects the client-specific gradient bias
  at each communication round.

\end{enumerate}

% ============================================================
\section{Thesis Contributions}
\label{sec:contributions}
% ============================================================

This thesis makes the following original contributions to the federated medical
imaging literature:

\begin{enumerate}

  \item[\textbf{C1.}] A three-stage centralised training pipeline --- combining
  Label-Distribution-Aware Margin loss, Sharpness-Aware Minimisation, Stochastic
  Weight Averaging, and classifier re-training (LDAM + ASAM + SWA + cRT) --- for
  breast DCE-MRI molecular subtype classification under severe class imbalance,
  evaluated over five-fold cross-validation with patient-level stratification.

  \item[\textbf{C2.}] A clinically motivated binary reformulation of the molecular
  subtype classification task that maps directly to the hormone therapy eligibility
  decision boundary, demonstrated to be more tractable under federated constraints
  than the four-class formulation.

  \item[\textbf{C3.}] An empirical comparison of FedAvg and SCAFFOLD aggregation
  strategies under three degrees of non-IID data heterogeneity
  (Dirichlet $\alpha \in \{0.1,\, 0.5,\, 100\}$) on the publicly available breast
  DCE-MRI dataset \cite{alves2023breastmri}, quantifying the convergence benefit
  of control-variate correction.

  \item[\textbf{C4.}] A proof-of-concept radiomic feature fusion experiment
  combining hand-crafted PyRadiomics features with deep MIL embeddings, providing
  an interpretable complement to the end-to-end deep learning pipeline.

\end{enumerate}
